\documentclass[14pt,a4paper]{extarticle}

\usepackage[T2A]{fontenc}
\usepackage[utf8]{inputenc}
\usepackage[russian]{babel}
\usepackage[a4paper,left=30mm,right=15mm,top=20mm,bottom=20mm]{geometry}
\usepackage{setspace}
\usepackage{indentfirst}
\usepackage{graphicx}
\usepackage{float}
\usepackage{hyperref}
\usepackage{xcolor}
\usepackage{listings}
\usepackage{longtable}
\usepackage{array}
\usepackage{tabularx}

\onehalfspacing
\setlength{\parindent}{1.25cm}

\hypersetup{
    colorlinks=true,
    linkcolor=blue,
    urlcolor=blue
}

\lstdefinestyle{sqlstyle}{
    language=SQL,
    basicstyle=\ttfamily\small,
    keywordstyle=\color{blue},
    commentstyle=\color{gray},
    stringstyle=\color{teal},
    breaklines=true,
    frame=single
}

\lstdefinestyle{pystyle}{
    language=Python,
    basicstyle=\ttfamily\small,
    keywordstyle=\color{blue},
    commentstyle=\color{gray},
    stringstyle=\color{teal},
    breaklines=true,
    frame=single
}

% ---- Данные для титульного листа (при необходимости меняй только здесь) ----
\newcommand{\studentfullname}{Пудовкина Татьяна Евгеньевна}
\newcommand{\studentgroup}{МО-\underline{\hspace{1.3cm}}}
\newcommand{\studentcourse}{\underline{\hspace{1.3cm}}}
\newcommand{\disciplinename}{SQL и получение данных}
\newcommand{\reporttheme}{Лабораторные работы №1 и №2}
\newcommand{\supervisorname}{Гулянов Руслан Гурамиевич}
\newcommand{\supervisorpost}{ассистент}
\newcommand{\directioncode}{02.03.03}
\newcommand{\directionname}{Математическое обеспечение и администрирование информационных систем}

\begin{document}

\begin{titlepage}
    \begin{center}
        \textbf{Министерство науки и высшего образования Российской Федерации}\\
        \textbf{ФГБОУ ВО «Омский государственный технический университет»}\\
        \vspace{1.2cm}
        Кафедра «Прикладная математика и фундаментальная информатика»\\
        \vspace{1.5cm}
        \textbf{\Large Лабораторная работа № 1, 2}\\
        \vspace{0.8cm}
    \end{center}

    \renewcommand{\arraystretch}{1.35}
    \begin{tabularx}{\textwidth}{@{}lX@{}}
        по дисциплине & \disciplinename \\
        на тему & \reporttheme \\
    \end{tabularx}

    \vspace{1.4cm}

    \begin{tabularx}{\textwidth}{@{}lX@{}}
        Студентка & \studentfullname \\
    \end{tabularx}
    {\footnotesize фамилия, имя, отчество полностью}

    \vspace{0.7cm}

    \begin{tabularx}{\textwidth}{@{}lXlX@{}}
        Курс & \studentcourse & Группа & \studentgroup \\
    \end{tabularx}

    \vspace{0.7cm}

    \begin{tabularx}{\textwidth}{@{}lX@{}}
        Направление & \directioncode\ \directionname \\
    \end{tabularx}
    {\footnotesize код, наименование}

    \vspace{0.9cm}

    \begin{tabularx}{\textwidth}{@{}lX@{}}
        Руководитель & \supervisorpost, \supervisorname \\
    \end{tabularx}
    {\footnotesize должность, ученая степень, звание; фамилия, инициалы}

    \vfill

    \begin{center}
        Омск, \the\year
    \end{center}
\end{titlepage}

\tableofcontents
\newpage

\section*{Введение}
\addcontentsline{toc}{section}{Введение}
В отчете представлены результаты выполнения первых двух лабораторных работ текущего года по дисциплине «SQL и получение данных»:
\begin{enumerate}
    \item ЛР1: Extract-Transform-Load (ETL).
    \item ЛР2: Разработка DWH и витрины данных.
\end{enumerate}

\section{Лабораторная работа №1. ETL}

\subsection{Цель работы}
Реализовать ETL-процесс: генерация синтетических данных, загрузка в PostgreSQL в неструктурированном виде, очистка аномалий средствами SQL и загрузка в структурированную таблицу.

\subsection{Постановка задачи}
\begin{enumerate}
    \item Сгенерировать синтетический датасет из 10 полей с намеренно добавленными аномалиями.
    \item Загрузить данные в таблицу \texttt{s\_sql\_dds.t\_sql\_source\_unstructured}.
    \item Реализовать SQL-функцию \texttt{fn\_etl\_data\_load(start\_date, end\_date)} для очистки и преобразования.
    \item Получить целевую таблицу \texttt{s\_sql\_dds.t\_sql\_source\_structured}.
    \item Реализовать верхнеуровневую функцию приложения \texttt{etl()} без входных параметров.
\end{enumerate}

\subsection{Структура проекта}
Вставьте дерево проекта (или скрин) после выполнения:
\begin{verbatim}
data-pipeline/
sql/
\end{verbatim}

\subsection{Описание реализации}
\subsubsection{Созданные таблицы PostgreSQL}
Кратко опишите назначение таблиц:
\begin{itemize}
    \item \texttt{t\_sql\_source\_unstructured} --- сырые данные (все поля VARCHAR).
    \item \texttt{t\_sql\_source\_structured} --- очищенные данные в корректных типах.
\end{itemize}

\subsubsection{SQL-функция очистки}
Вставьте код или ключевой фрагмент:
\begin{lstlisting}[style=sqlstyle,caption={Фрагмент fn_etl_data_load}]
-- вставить ваш SQL-код
\end{lstlisting}

\subsubsection{Код приложения}
\begin{lstlisting}[style=pystyle,caption={Верхнеуровневая функция etl()}]
# вставить/описать вызовы:
# get_dataset()
# load_data_to_db()
# fill_structured_table()
\end{lstlisting}

\subsection{Проверка результатов}
Заполните таблицу фактическими значениями:
\begin{longtable}{|p{0.55\textwidth}|p{0.35\textwidth}|}
\hline
Проверка & Результат \\ \hline
Количество строк в t\_sql\_source\_unstructured & \\ \hline
Количество строк в t\_sql\_source\_structured & \\ \hline
Минимальная/максимальная дата в целевой таблице & \\ \hline
\end{longtable}

\subsection{Вывод по ЛР1}
Краткий вывод: что реализовано, какие аномалии обрабатываются, как подтверждена работоспособность.

\newpage
\section{Лабораторная работа №2. DWH и витрина данных}

\subsection{Цель работы}
Построить модель DWH (звезда/снежинка), загрузить таблицы-справочники, сформировать таблицу витрины, создать представление и реализовать перекладку данных в MySQL.

\subsection{Постановка задачи}
\begin{enumerate}
    \item Создать таблицы-справочники с полями \texttt{id}, \texttt{name}.
    \item Создать таблицу \texttt{t\_dm\_task} с ID справочников и фактами.
    \item Реализовать функцию \texttt{fn\_dm\_data\_load(start\_dt, end\_dt)}.
    \item Создать витрину \texttt{v\_dm\_task} с явным перечислением полей.
    \item Создать таблицу в MySQL и выполнить перенос данных из PostgreSQL.
\end{enumerate}

\subsection{Схема данных}
Вставьте ER-диаграмму (скрин из DBeaver/Draw.io/Figma):
\begin{figure}[H]
    \centering
    % \includegraphics[width=0.95\textwidth]{images/dwh_schema.png}
    \caption{Схема данных (звезда/снежинка)}
\end{figure}

\subsection{Описание реализации}
\subsubsection{Справочники и факт}
\begin{itemize}
    \item \texttt{t\_dim\_customer}
    \item \texttt{t\_dim\_city}
    \item \texttt{t\_dim\_product\_category}
    \item \texttt{t\_dim\_status}
    \item \texttt{t\_dim\_source\_system}
    \item \texttt{t\_dm\_task}
\end{itemize}

\subsubsection{Функция заполнения витрины}
\begin{lstlisting}[style=sqlstyle,caption={Фрагмент fn_dm_data_load}]
-- вставить ваш SQL-код
\end{lstlisting}

\subsubsection{Представление v_dm_task}
\begin{lstlisting}[style=sqlstyle,caption={Создание v_dm_task}]
-- перечислить поля явно, без SELECT *
\end{lstlisting}

\subsubsection{Перекладка данных в MySQL}
\begin{lstlisting}[style=pystyle,caption={Функция переноса в MySQL}]
# transfer_dm_to_mysql(...)
\end{lstlisting}

\subsection{Проверка результатов}
\begin{longtable}{|p{0.55\textwidth}|p{0.35\textwidth}|}
\hline
Проверка & Результат \\ \hline
Количество строк в PostgreSQL: v\_dm\_task & \\ \hline
Количество строк в MySQL: t\_dm\_task & \\ \hline
Корректность ключевых полей (ID, факт) & \\ \hline
\end{longtable}

\subsection{Вывод по ЛР2}
Краткий вывод: какая схема построена, как формируется витрина, как подтвержден перенос в MySQL.

\newpage
\section*{Заключение}
\addcontentsline{toc}{section}{Заключение}
В ходе выполнения ЛР1 и ЛР2 реализован полный контур обработки данных:
\begin{enumerate}
    \item ETL в PostgreSQL с очисткой аномалий.
    \item Формирование DWH-структуры и витрины данных.
    \item Перекладка витрины в MySQL из приложения.
\end{enumerate}

\section*{Приложение А (опционально)}
\addcontentsline{toc}{section}{Приложение А (опционально)}
Сюда можно вынести полный SQL-код и длинные листинги Python.

\end{document}
